In \textit{I Flew Over Navajo Land}, Tim Yazzie recounts his experience as part of a selected group tasked with delivering aid by plane to Navajo communities affected by severe winter weather. The journey began on February 2 in 1949, as Yazzie and three other men set out from Tuba City, eventually boarding separate planes in Phoenix to survey and supply Navajo communities. Yazzie describes the flight over significant villages and sites such as Shonto, Oraibi, and Mexican Hat, and his role as a guide familiar with the terrain. Throughout the flights, they dropped essential supplies—hay, food, and firewood—over villages, witnessing people hurriedly gathering the resources on the ground. Yazzie describes the harsh landscape from above, observing cattle and horses immobilized by crusted snow, suffering animals, and sheep corralled for survival.

In the days that followed, Yazzie and his group met with officials, including the Governor of Arizona and an army officer, to discuss the challenges faced by the Navajo people. On multiple flights, Yazzie's planes encountered intense fog, turbulent weather, and even a motor fire, depicting the physical dangers faced by those delivering aid. This narrative provides insight into both the logistical struggles of providing aid by air and the broader emotional toll on Navajo communities amid environmental adversities.